\documentclass[10pt, a4paper]{article}

\usepackage[utf8]{inputenc}
\usepackage[T1]{fontenc}
% headsep keeps the date in the header close to the name, well clear of the page edge
\usepackage[left=2.0cm, right=2.0cm, top=1.9cm, bottom=1.6cm, headsep=0.5cm]{geometry}
\usepackage{titlesec}
\usepackage{xcolor}
\usepackage{fontawesome5}
\usepackage{tikz}
\usepackage{xparse}
\usepackage{amsfonts}
\usepackage{amsmath}
\usepackage{fancyhdr}
\usepackage[style=iso]{datetime2}
\usepackage{url}
\usepackage{enumitem}
\usepackage[hidelinks]{hyperref} % clickable links; load last

\setlength{\headheight}{13.6pt}

\usetikzlibrary{shapes.geometric}

\renewcommand{\familydefault}{\rmdefault}

% Product names are never split across lines
\hyphenation{Antigravity ComfyUI ChatGPT Likeligood Robostrat Pentech}

% Define a custom command for the hexagonal icon
\newcommand{\hexicon}[2]{%
  \smash{% Tells LaTeX to ignore the vertical height of the icon
    \begin{tikzpicture}[baseline={([yshift=-0.6ex]char.center)}]
    \node[
      draw=#1,
      fill=white,
      shape=regular polygon,
      regular polygon sides=6,
      inner sep=-1.0pt,
      minimum size=0.75cm,
      line width=1.4pt
    ] (char) {\color{#1}\normalsize#2};
    \end{tikzpicture}%
  }%
}

\definecolor{cvDarkgray}{gray}{0.2}    % Define colors for the classic grey theme

\definecolor{cvGreen}{HTML}{006400}    % Dark Green
\definecolor{cvBlue}{HTML}{1A5276}     % Steel Blue: Strong, corporate, and clean.
\definecolor{cvTeal}{HTML}{0E6251}     % Deep Sea: Modern, sophisticated, and easy on the eyes.
\definecolor{cvMidnight}{HTML}{212F3D} % Midnight Blue: Very dark, looks almost black until compared.

\definecolor{cvPurple}{HTML}{4A235A}   % Royal Plum: Excellent contrast with cvDarkgray text.
\definecolor{cvBurgundy}{HTML}{641E16} % Oxblood: Serious and authoritative, prints very well.
\definecolor{cvIndigo}{HTML}{283593}   % Academic Indigo: Slightly more vibrant than Steel Blue.
\definecolor{cvBronze}{HTML}{784212}   % Burned Umber: Distinctive, warm, and professional.

% Every hexagon sits centered in a column of \hexcolw, and its text starts \hexgap after that column,
% so the contact rows line up with the section headings.
\newlength{\hexcolw}\setlength{\hexcolw}{1.0cm}
\newlength{\hexgap}\setlength{\hexgap}{0.24cm}

% Custom CV Section Command
% Usage: \cvsection{Color}{Icon}{Title}
\NewDocumentCommand{\cvsection}{ m m m }{%
  \par% The \nobreak below keeps a heading on the same page as its first entry
  \vspace{2.2ex}
  \noindent
  \begin{minipage}[c]{\hexcolw}
    \centering \hexicon{#1}{#2}
  \end{minipage}%
  \begin{minipage}[c]{\dimexpr\linewidth-\hexcolw\relax}
    {\Large\bfseries\color{cvDarkgray}\hspace{\hexgap}#3}\par
    \vspace{-1.6ex}
    \color{#1!70}\rule{\linewidth}{1.5pt}
  \end{minipage}%
  \par\nobreak\vspace{1.2ex}% No page break between a heading and its first entry
}

% Custom CV entry command
\newcommand{\cventry}[5]{%
  \noindent
  \begin{minipage}[t]{3.3cm}
  \vspace{0pt}\textcolor{cvDarkgray}{\textbf{#1}}
  \end{minipage}%
  \begin{minipage}[t]{\dimexpr\linewidth-3.3cm\relax}
  \vspace{0pt}\textbf{#2}\ifx&#3&\else, #3\fi\ifx&#4&\else, \textit{#4}\fi\par
  \ifx&#5&\else \vspace{0.08cm} #5 \fi
  \end{minipage}\par\vspace{0.26cm}
}

% Custom CV item command
\newcommand{\cvitem}[2]{%
  \noindent
  \begin{minipage}[t]{3.3cm}
  \vspace{0pt}\textcolor{cvDarkgray}{\textbf{#1}}
  \end{minipage}%
  \begin{minipage}[t]{\dimexpr\linewidth-3.3cm\relax}
  \vspace{0pt}#2
  \end{minipage}\par\vspace{0.16cm}
}

% Header and Footer setup
\pagestyle{fancy}
\fancyhf{}
\renewcommand{\headrulewidth}{0pt}
\rhead{\textcolor{lightgray}{\small \today}}

% Contact header: two columns, each an icon column beside a text column.
% Usage: \contactcolumn{Width}{Icons}{Texts}, one row per icon, rows separated by \\[0.3cm]
% The width is fixed in \contactw first: inside the outer minipage \linewidth is already the column's.
\newlength{\contactw}
\newcommand{\contactcolumn}[3]{%
  \setlength{\contactw}{#1}%
  \begin{minipage}[t]{\contactw}%
    \begin{minipage}[t]{\hexcolw}\centering#2\end{minipage}%
    \hspace{\hexgap}%
    \begin{minipage}[t]{\dimexpr\contactw-\hexcolw-\hexgap\relax}#3\end{minipage}%
  \end{minipage}%
}

% The public contact details. A private variant is built locally by defining \privatecontact as the
% path of a file that redefines these macros (and fills the right-hand column); that file is never
% part of this repository.
\newcommand{\cvemail}{\href{mailto:oscar@anghammar.com}{\texttt{oscar@anghammar.com}}}
\newcommand{\cvrighticons}{}
\newcommand{\cvrighttexts}{}
\ifdefined\privatecontact \input{\privatecontact}\fi

\begin{document} % -----------------------------------------------------------------------------------------------------

% Header Information
  \noindent{\Huge\bfseries\textcolor{cvDarkgray}{Oscar Anghammar}}\\[0.45cm]
  \noindent
  \contactcolumn{0.45\linewidth}{%
    \hexicon{cvDarkgray}{\faEnvelope}\\[0.3cm]
    \hexicon{cvDarkgray}{\faDesktop}\\[0.3cm]
    \hexicon{cvDarkgray}{\faGithub}\\[0.3cm]
    \hexicon{cvDarkgray}{\LARGE\textbf{\$}}%
  }{%
    \cvemail\\[0.3cm]
    \url{https://anghammar.com}\\[0.3cm]
    \url{https://github.com/Anghammar}\\[0.3cm]
    \textbf{Anghammar konsult AB}\\
    Org.\ no.\ 559509-3542%
  }%
  \contactcolumn{0.55\linewidth}{\cvrighticons}{\cvrighttexts}%
  \par\vspace{0.3cm}

  \cvsection{cvPurple}{\large\textbf{$\genfrac{}{}{0pt}{2.0pt}{\text{TL;}}{\text{DR}}$}}{Summary}
  \noindent Software engineer with a Master's degree in mathematics and nine years of professional experience
  in fintech, integrating mathematics with production code. My past work includes designing asset allocation
  algorithms, end-to-end probability models for automated trading, and calibration algorithms together with
  live market data pipelines for automatic model parameter updates. Since 2026 running my own company,
  Anghammar konsult AB, working on reinforcement learning, LLM agents, local AI deployment and
  serverless mobile apps.
  \vspace{0.15cm}

  \cvsection{cvMidnight}{\faBriefcase}{Experience}
  \cventry{2026--}{Founder / Software Engineer}{Anghammar konsult AB}{Stockholm}{
    Software development and consulting; see Projects for my recent work.
  }
  \cventry{2018--2026}{Mathematician / Software Engineer}{Likeligood AB}{Stockholm}{
    Fintech startup (4--5 developers) developing a robot for semi-automated trading, mainly in sports betting.
    I was with the company from its start until it ceased operations.
    \begin{itemize}[topsep=3pt, partopsep=0pt, itemsep=1pt, parsep=0pt]
      \item \textbf{Built end-to-end probability models} for live soccer betting, from scraping data on thousands
        of matches to modeling, testing and performance analysis: Markov chains calibrated live against market
        data with an Unscented Kalman Filter, implemented as Monte Carlo simulations on the GPU in GLSL and
        Vulkan.
      \item \textbf{Generalized the modeling framework} into an API designed to fit most sports.
      \item \textbf{Designed and deployed} an automated cryptocurrency trading robot.
      \item \textbf{Integrated our solution into a partner's codebase} in a joint venture project with a major
        sports data provider, working closely with their engineers.
      \item \textbf{Ran live operations} on our Linux servers: technical contact for the system operators, and
        debugging live failures from server logs.
      \item \textbf{Contributed to the UI} by designing the parts for controlling and monitoring live data
        filtering and model parameters.
      \item \textbf{Led the recruitment and onboarding} of a junior engineer.
    \end{itemize}
    \textbf{Technologies:} Scala (main production language), MATLAB, JavaScript, GLSL, Circe,
    Apache Commons Math and Protobuf.
  }
  \cventry{2017--2018}{Mathematician}{Pentech Ltd}{Remote}{
    Developed and implemented asset allocation algorithms for pension funds.
    Analyzed financial data and built prototypes in MATLAB, with production implementation in C\#.
  }
  \cventry{2011--2017}{Teaching Assistant}{Chalmers University of Technology}{Gothenburg}{
    Led exercise sessions for undergraduate courses, mainly in calculus and linear algebra, and MATLAB
    computer exercises; mentored students in a preparatory mathematics course.
  }

  \cvsection{cvIndigo}{\LARGE$\mathbb{R}$}{Technical Proficiencies}
  \cvitem{Generative AI}{
    Tool-using LLM agents (MCP); agentic development (Claude Code, Antigravity); self-hosted LLM/VLM
    serving (llama.cpp) behind Cloudflare; image generation (ComfyUI).
  }
  \cvitem{ML}{
    Self-play reinforcement learning (PPO) with PyTorch on ROCm; ONNX in the browser.
  }
  \cvitem{Mathematics}{
    Strong foundation in probability modeling, including Markov chains, Monte Carlo simulation and Kalman
    filtering. Extensive academic background in linear and multilinear algebra, abstract algebra, real and
    complex analysis, and topology.
  }
  \cvitem{Tools, OS}{Git, IntelliJ, SBT, \LaTeX; Windows, macOS and Linux, all three in regular use.}

  \cvsection{cvGreen}{\faEdit}{Projects}
  \cvitem{Robostrat}{
    A turn-based strategy game developed alongside its own AI, as a testbed for an AI training workflow
    intended to carry over to other turn-based strategy games. The bot is trained by self-play reinforcement
    learning (PyTorch on ROCm, about a thousand games at a time) and plays in the browser via ONNX.
    For efficient training, the rules first ran as a PyTorch copy on the GPU, and have since been migrated to a Rust core.
    The training engine is kept in step with the browser's TypeScript engine by differential fuzzing.
    Language-model scouts (mostly Gemini) play the bot through an MCP server and report the weaknesses they find. 
    Their view of the board is text that does the hex geometry for them (reach, supply lines, attacks priced by playing them out),
    leaving the strategy to the model.
    \\[0.06cm]
    \url{https://anghammar.com/robostrat/}
  }
  \cvitem{Wiki Harness\\Racing}{
    A zero-backend browser harness that benchmarks LLMs as tool-using agents on automated Wikiraces:
    self-correction from tool-call errors, loop detection and live speed metrics. It runs against any
    OpenAI-compatible endpoint, such as a local llama.cpp server; the public demo uses an LLM hosted on my
    own server.
    \\[0.06cm]
    \url{https://harness.anghammar.com}
  }
  \cvitem{Pröjsa}{
    Splits restaurant bills with the user's own AI assistant (Gemini, Claude or ChatGPT), which
    turns the receipt into a link; everyone taps their items and pays with Swish. Privacy by design:
    Pröjsa never sees images and has no accounts or tracking.
    \\[0.06cm]
    \url{https://projsa.anghammar.com}
  }
  \cvitem{KnotNet}{
    An experimental serverless messaging app for iOS and Android, in public beta. Phones sync posts and
    messages directly over Bluetooth LE and relay them for each other (store-and-forward), with no servers
    and no internet. Native apps in Swift and Kotlin, with a gossip sync protocol, signed identities and
    end-to-end encrypted DMs. Developed and tested with a hardware-in-the-loop harness that drives real
    phones over USB.
    \\[0.06cm]
    \url{https://anghammar.com/knotnet/}
  }

  \cvsection{cvBurgundy}{\faCode}{Programming}
  \cvitem{Scala}{\textbf{Production systems.}
    Eight years as my main language at Likeligood, in a mixed functional and object-oriented style.
    Currently my preferred programming language.
  }
  \cvitem{Python}{\textbf{ML and generative AI.}
    Robostrat's training pipeline (PyTorch), test suite and MCP server; custom and patched ComfyUI nodes
    for AMD GPUs; everyday scripting.
  }
  \cvitem{TypeScript/\\JavaScript}{\textbf{Web apps.}
    Data scrapers and UI fixes at Likeligood; Robostrat's game engine and browser client, and Wiki
    Harness Racing.
  }
  \cvitem{Java}{\textbf{The JVM ecosystem.}
    Integrating Java libraries into Scala code and debugging enterprise Java codebases at Likeligood.
    I also read Kotlin, which is close to Scala.
  }
  \cvitem{MATLAB}{\textbf{Numerical computing.}
    Used extensively since 2011 (university, Pentech, Likeligood), mainly numerical work and linear
    algebra. My preferred language before Scala.
  }
  \cvitem{C\#}{\textbf{Production algorithms.}
    Ten months implementing asset allocation algorithms in production at Pentech.
  }

  \cvsection{cvBlue}{\faGraduationCap}{Education}
  \cvitem{2015--2017}{Various university courses; mainly graduate-level courses in mathematics.}
  \cventry{2013--2015}{Master of Science in Mathematics}{University of Gothenburg}{}{
    \textbf{Master's Thesis:} \emph{de Rham Cohomology for Grassmann and Flag Manifolds} (Supervisor: Genkai Zhang).
    Studies singular and de Rham cohomologies for the Grassmannian, exploring Schubert calculus
    and a possible generalization of Horn's conjecture.
  }
  \cventry{2010--2013}{Bachelor of Science in Mathematics}{University of Gothenburg}{}{}

  \cvsection{cvBronze}{\faInfo}{Other}
  \cvitem{Languages}{Swedish (native), English (fluent).}
  \cvitem{Driving license}{B (car).}

\end{document} % -------------------------------------------------------------------------------------------------------
